% Chapter 12: loading new firmware.
\chapter{Loading new firmware (flashing)}\label{ch:flashing}
\index{firmware!updating}\index{flashing}

The program inside the controller is called the \emph{firmware}. This manual
is for version fw=\fwver. You only need to update it if your supplier sends a
new version. Updating is a technician job; it needs a computer with the source
repository and its tools set up (see \texttt{docs/runbooks/flash.md} in the
repository).

\begin{warning}
Switch the motor power off, or at least make sure nobody is near the roll,
while you flash. The relay is opened while the controller restarts.
\end{warning}

\section{Steps}
\begin{enumerate}
  \item Connect the computer to the controller's \textbf{own USB port} (not the
        serial adapter).
  \item In the repository folder run:\par
        \texttt{uv run python scripts/flash.py}
  \item The script builds the firmware, then prints\par
        \shows{Double-tap RST on the Pro Micro now.}\par
        Quickly press the controller board's reset button twice. The board
        then waits about 8 seconds for the new firmware.
  \item The script finds the board and loads the firmware. When it reports
        success the controller restarts and shows the idle screen.
\end{enumerate}

On Windows with WSL, a one-time setup is needed first:
\texttt{uv run python scripts/flash.py --setup} (installs \texttt{usbipd} and
asks for administrator permission once).

\section{If flashing fails}
\begin{center}
\renewcommand{\arraystretch}{1.4}
\begin{tabularx}{\textwidth}{@{}>{\raggedright\arraybackslash}p{0.27\textwidth}>{\raggedright\arraybackslash}X@{}}
\toprule
\textbf{Message} & \textbf{What to do} \\
\midrule
\texttt{AttachTimeout} & The board was not handed to WSL in time. Double-tap
  reset again sooner after the prompt; check the USB cable. \\
\texttt{PortTimeout} & The board's port did not appear. Double-tap again; try
  another USB cable or port (some cables carry power only). \\
\texttt{AvrdudeFailed} & Loading failed part way. Run the script again. The
  board can always be re-flashed with a double-tap. \\
\bottomrule
\end{tabularx}
\end{center}

\section{After flashing}
\begin{itemize}
  \item If you came from a version older than fw=6, the machine starts on
        factory settings. Calibrate and save once (Chapter~\ref{ch:calibration}).
  \item Run the connection check (Section~\ref{sec:doctor}); it confirms the
        firmware is compatible.
  \item Wind one roll and check the result.
\end{itemize}
